% tab:war
\begin{table}[htbp]
\centering\small
\caption{Demographic and economic impact of the two world wars (percentage change)}
\label{tab:war}
\begin{threeparttable}
\begin{tabular}{lrrr}
\toprule
& Population & GDP & GDP per capita \\
\midrule
\multicolumn{4}{l}{\textit{(a) First World War: decline from 1913 to the 1918/19 trough}} \\
Austria & $-$5 & $-$38 & $-$35 \\
Belgium & 0 & $-$32 & $-$32 \\
France & $-$7 & $-$25 & $-$31 \\
USSR & $-$1 & $-$60 & $-$59 \\
Lithuania & $-$36 & $-$60 & $-$38 \\
Latvia & $-$35 & $-$59 & $-$37 \\
Estonia & $-$9 & $-$29 & $-$22 \\
\midrule
\multicolumn{4}{l}{\textit{(b) Second World War: decline from 1938 to the 1944--6 trough}} \\
Greece & 3 & $-$64 & $-$65 \\
Germany & $-$6 & $-$58 & $-$56 \\
Austria & 1 & $-$51 & $-$52 \\
Netherlands & 6 & $-$47 & $-$50 \\
Poland & $-$23 & $-$57 & $-$49 \\
Lithuania & $-$5 & $-$48 & $-$48 \\
Latvia & $-$23 & $-$40 & $-$28 \\
Estonia & $-$20 & $-$42 & $-$29 \\
\bottomrule
\end{tabular}
\begin{tablenotes}[para,flushleft]
\footnotesize
\textit{Notes:} each column is measured from the pre-war year to its own trough within the window named. Lithuania in panel (b) within actual borders (the 1939 loss of Klaip\.{e}da and gain of Vilnius); within post-war borders its population fell by 22 per cent. \textit{Sources:} Baltic states: this paper and the companion series; others: \textcite{mpd2023}.
\end{tablenotes}
\end{threeparttable}
\end{table}
